\chapter{\introductionname}

El olfato humano es uno de los sentidos químicos fundamentales y desempeña un papel clave en la detección del entorno, la seguridad alimentaria, la memoria y la emoción. A diferencia de la visión y la audición, que procesan estímulos físicos (luz y sonido), el olfato traduce señales químicas volátiles presentes en el aire en experiencias perceptivas complejas. Esta capacidad permite identificar alimentos en buen o mal estado, reconocer peligros (como humo o gases), y desencadenar respuestas afectivas y evocaciones memorísticas profundas~\cite{BuckAxel1991,Kandel2013}.

Para replicar esta capacidad sensorial, se han desarrollado dispositivos conocidos como narices electrónicas (e-noses). Una nariz electrónica es un dispositivo que se considera parte de la rama de investigación llamada biomímesis, ya que imita el sentido del olfato humano mediante el uso de sensores químicos para detectar y reconocer olores y compuestos volátiles en el aire. Estos dispositivos están diseñados para identificar patrones específicos de gases y vapores, permitiendo su aplicación en diversas áreas como la seguridad, la industria alimentaria, la medicina y el control ambiental.

El olfato canino destaca por su extraordinaria sensibilidad y discriminación. En comparación con los humanos, los perros poseen un epitelio olfatorio mucho más extenso y un número de receptores olfatorios significativamente mayor, lo que amplifica su capacidad para detectar concentraciones ínfimas de compuestos volátiles. Además, su morfología nasal favorece una mayor captura y análisis de olores. Estas características, combinadas con entrenamiento y aprendizaje asociativo, permiten a los perros distinguir mezclas complejas y realizar tareas de detección en ámbitos como seguridad, búsqueda y rescate, medicina y control ambiental.

Por ende, este proyecto plantea estudiar la viabilidad de una nariz electrónica como herramienta de apoyo para los agentes caninos en tareas de detección de sustancias. Se explorará si la tecnología de narices electrónicas puede complementar y potenciar las capacidades olfativas de los perros, mejorando la eficiencia y precisión en la identificación de sustancias específicas en diversos contextos operativos.

\section{Motivación}
Desde una perspectiva personal, este trabajo nace de una profunda admiración por los cuerpos de seguridad y, en particular, por las unidades caninas. %Su profesionalidad, disciplina y entrega representan valores que inspiran la intención de unirme a este cuerpo en el futuro y contribuir activamente a su misión.

Mi vínculo con los animales, especialmente con los perros, es cercano y afectivo. He convivido con ellos desde pequeño gracias a que mi abuelo siempre ha tenido varios perros en casa. Esa experiencia cotidiana, hecha de cuidado, respeto y aprendizaje, ha reforzado tanto mi cariño hacia los animales como mi interés por comprender y aprovechar científicamente sus capacidades olfativas en beneficio de la sociedad.

La motivación que guía este proyecto es, por tanto, doble: por un lado, aportar soluciones tecnológicas que apoyen el desempeño de las unidades caninas en tareas de seguridad, rescate y detección; por otro, honrar el papel de los perros como compañeros y servidores en labores críticas, integrando sensibilidad humana y rigor científico para lograr un impacto positivo.
